\documentclass[english,10pt,a4paper,onecolumn]{article}
\usepackage[T1]{fontenc}
\usepackage[width=17cm]{geometry}
\usepackage{babel}
\usepackage{amsmath}
\usepackage{enumitem}
\title{Theory-Native Feature Isolation}
\author{Feature Isolation Project}
\date{}
\begin{document}
\maketitle

\section{Scope}
This note documents the feature-isolation model implemented by the current
project state. It assumes the theory-native SPL DSL, a valid-system space
computed from the conjunction of asserted formulas, and a theory feature
registry derived from that system space.

\section{Inputs}
Feature isolation operates on:
\begin{itemize}[leftmargin=*]
\item the valid system space $S_{valid}$,
\item the available system space $S_{available}$,
\item the missing system space $S_{missing}$,
\item and the theory feature registry generated from $S_{valid}$.
\end{itemize}

Only available valid systems participate in the active witness search.
Missing systems influence interpretation: a witness may isolate more than one
candidate feature.

\section{Partitioning}
For a feature $t$, the implementation partitions the available valid systems
into:
\[
I(t) = \{ s \in S_{available} \mid t \text{ is present in } s \}
\]
\[
U(t) = \{ s \in S_{available} \mid t \text{ is absent in } s \}
\]

A witness is a pair $(I, U)$ with $I \subseteq I(t)$ and $U \subseteq U(t)$.

\section{Flattening}
Each system identifier is mapped to the set of feature identifiers present in
that system. The implementation uses two flattening operators:
\[
\widehat{I} = \bigcap_{s \in I} F(s)
\]
\[
\widetilde{U} = \bigcup_{s \in U} F(s)
\]
where $F(s)$ denotes the feature IDs present in system $s$.

The isolated candidate set of a witness is:
\[
\overline{D}(I, U) = \widehat{I} \setminus \widetilde{U}
\]

\section{Interpretation of results}
If $|\overline{D}(I, U)| = 1$, the witness is unique with respect to the
available information. If $|\overline{D}(I, U)| > 1$, the witness is
ambiguous: it isolates a class of candidate features instead of a singleton.

This ambiguity is expected in the presence of missing systems. The project does
not collapse such cases to an arbitrary single feature.

\section{Minimal witnesses}
The implementation searches for minimal witnesses that isolate a fixed target
candidate set. A witness $(I, U)$ is minimal if no proper reduction of $I$ or
$U$ isolates the same candidate set.

The result reported for a feature therefore consists of:
\begin{itemize}[leftmargin=*]
\item the isolated candidate feature set, and
\item the set of minimal witness pairs that isolate exactly that set.
\end{itemize}

\section{Reporting}
The visible report labels each result as either:
\begin{itemize}[leftmargin=*]
\item \texttt{UNIQUE}, if the candidate feature set has cardinality 1, or
\item \texttt{AMBIGUOUS}, otherwise.
\end{itemize}

The report also prints the witness count and the candidate feature IDs.

\section{Limitations}
The current implementation is theory-correct but exhaustive: it constructs the
full feature registry derived from the valid-system space. For larger SPLs, the
main limit is therefore combinatorial growth of the feature space rather than
an artificial fixed bound inside the witness search.

\end{document}
